% Appendix F, the construction pilot (restructure plan, part 3). Every number is a macro of generated/pilot/numbers.tex
% (\plv; evidence_ledger_pilot.md: the frozen run results/runs/pilot-v0.3.3, its comparison files, the B1 check-origin
% files and verifier specs, and the exploratory re-analysis of D49), of the design constants (\dv) or of the fresh
% numbers (\fv), or a cell of a generated table (the pilot's tables, frozen at c5ede4c: evidence_ledger.md; the
% exploratory tables: evidence_ledger_exploratory.md). Moved here: results_pilot.tex but what Section sec:pilot keeps,
% analysis.tex's pilot cases, protocol.tex's held-out pool, limitations.tex's pilot items, and the outcome matrix.
\section{The construction pilot}
\label{app:pilot}

\subsection{Run, pool, labels and conventions}
\label{sec:results-pilot}\label{sec:pools}
Every rate here comes from run \texttt{\PILOTRUN}: \plv{run.families.w} task families, \plv{run.conditions.w} conditions
and \plv{run.candidates.w} candidates per family, giving \plv{run.pairs} graded pairs, every pair graded under harness
\plv{run.harness}, all at the same code revision, with no infrastructure outcome (\plv{run.infrastructure}). Since protocol
v0.7 the pilot is a development study, and it shares its families and verifiers with \texttt{\FRESHRUN}, so the runs are
not independent evidence (Section~\ref{sec:pilot}). Counts about verifier construction come from the verifier specs
(Table~\ref{tab:verifiers}); other counts come from the decision log, as cited.

Per family, the pool holds the agent's own pre-correction state; \dv{pilot.pool.alternatives.w} independently written
valid alternatives, implemented from the reference instruction without the reference diff; and
\dv{pilot.pool.nearmisses.w} near misses, each built from the reference to violate a targeted requirement (the rule's
labels record from \plv{nm.multi.min} to \plv{nm.multi.max} violated requirements for \plv{nm.multi.n.w} of them).
Such near misses may be easier or harder to catch than the errors that agents make on their own, which is what the fresh
attempts add. The no-op and the reference are graded as sanity checks only. Alternatives and near misses were generated
after every evaluated verifier of their family was frozen (the rebuild of a B1 is in Appendix~\ref{app:construction}),
by separate agents told not to read the frozen verifiers or registered candidates, which were readable; the audits of
their transcripts and the other departures from blindness are in Appendix~\ref{app:deviations}.

Labels are the proposals of rule D19 (Table~\ref{tab:label-rules}), not human labels, so this appendix is descriptive.
Of the \plv{heldout} held-out candidates, \plv{labels.invalid} are labeled invalid, \plv{labels.valid} valid and
\plv{labels.ambiguous} ambiguous; the ambiguous ones are left out of the rates (Table~\ref{tab:labels}). The tables are
generated by \texttt{paper/make\_tables.py}, every number in them is listed in \texttt{paper/evidence\_ledger.md}, and these
conventions apply to them:
\begin{itemize}
\item \emph{Degenerate intervals.} With \plv{run.families.w} families, a cluster-bootstrap interval is degenerate (of no
  width, marked $^{*}$ in every table) whenever every family has the same rate or the same difference. For rates the
  tables therefore also print the exact binomial interval, which ignores clustering and serves only as a reference scale;
  no exact interval is given for paired differences. For example, for \plv{c3.fa} the upper end of the exact
  \dv{design.exact.level}\% interval is \plv{c3.fa.exact.upper}.
\item \emph{Exploratory analyses.} Analyses marked exploratory were chosen after results had been seen and are never
  primary: the construction-intent relabeling (Table~\ref{tab:sensitivity}) in protocol v0.3, after the superseded runs
  \texttt{pilot-v0.2} and \texttt{pilot-v0.2-h2}; the post hoc analyses (Table~\ref{tab:posthoc}) in protocol v0.4,
  after the internal review of the first final run, with the duplicates view widened in v0.5.
\item \emph{Frozen names.} These tables are frozen with the pilot's run (check~e of
  \path{scripts/verify_frozen_inputs.py}) and keep the names they were frozen with: ``C1 initial'' and ``C2 initial'' are the initial-only verifiers,
  ``C3 full'' and ``C4 full'' the correction-aware ones, and ``B2 full-session LLM'' is the session-LLM baseline.
  ``Reward fidelity'' in the caption of Table~\ref{tab:summary} means false acceptance and false rejection on the
  constructed pool and nothing more: no verifier was used as a reward and no training was run. ``Pre-specified'' in
  their captions means planned in the locally logged protocol before any evaluation run (protocol~\S7, v0.1;
  Appendix~\ref{app:protocol}).
\end{itemize}

\IfFileExists{generated/results_summary.tex}{\begin{table}[t]\centering\footnotesize
\caption{Verifier reward fidelity on the held-out pool, run \texttt{pilot-v0.3} (4 tasks, 4 families). \textbf{Automated pilot: labels are rule-based proposals (D19), not human labels.} Rates with (num/den) and 95\% task-family bootstrap interval. Reference solution and no-op excluded; ambiguous candidates excluded. $J$ = (1$-$FR)$-$FA.}\label{tab:summary}
\begin{tabular}{lllrrr}\toprule Condition & False acceptance & False rejection & $J$ & Tasks w/ accepted neg. & Infra \\\midrule
C1 initial, plain & 0.35 (6/17) [0.07, 0.71] & 0.20 (2/10) [0.00, 0.67] & +0.45 & 3/4 & 0/32 \\
C2 initial, audited & 0.29 (5/17) [0.00, 0.63] & 0.20 (2/10) [0.00, 0.67] & +0.51 & 2/4 & 0/32 \\
C3 full, plain & 0.00 (0/17) [0.00, 0.00] & 0.00 (0/10) [0.00, 0.00] & +1.00 & 0/4 & 0/32 \\
C4 full, audited & 0.00 (0/17) [0.00, 0.00] & 0.00 (0/10) [0.00, 0.00] & +1.00 & 0/4 & 0/32 \\
B2 full-session LLM & 0.06 (1/17) [0.00, 0.14] & 0.00 (0/10) [0.00, 0.00] & +0.94 & 1/4 & 0/32 \\
B1 repository tests & 0.29 (5/17) [0.11, 0.64] & 0.40 (4/10) [0.00, 0.82] & +0.31 & 3/4 & 0/32 \\
\bottomrule\end{tabular}\end{table}
}{\unfilled{results table}}
\IfFileExists{generated/negatives_table.tex}{\begin{table}[t]\centering\footnotesize
\caption{Accepted invalid candidates by origin and by negative type (accepted/total) (automated pilot). Correction-violating: violates a reference requirement introduced after the first user request.}\label{tab:negatives}
\begin{tabular}{lrrrr}\toprule Condition & Pre-correction state & Near miss & Correction-violating & Initial-violating \\\midrule
C1 initial, plain & 2/2 & 4/15 & 6/16 & 0/1 \\
C2 initial, audited & 2/2 & 3/15 & 5/16 & 0/1 \\
C3 full, plain & 0/2 & 0/15 & 0/16 & 0/1 \\
C4 full, audited & 0/2 & 0/15 & 0/16 & 0/1 \\
B2 full-session LLM & 0/2 & 1/15 & 1/16 & 0/1 \\
B1 repository tests & 0/2 & 5/15 & 5/16 & 0/1 \\
\bottomrule\end{tabular}\end{table}
}{\unfilled{negatives table}}

\subsection{The pilot's rates}
\paragraph{Initial-only against correction-aware verifiers.} C1 accepts \plv{c1.fa} invalid held-out candidates and
rejects \plv{c1.fr} valid ones; the same generation procedure given the correction-aware contract (C3) accepts
\plv{c3.fa} and rejects \plv{c3.fr} (Tables~\ref{tab:summary} and~\ref{tab:negatives}). The paired differences are
\plv{cmp.c1-c3.fa.diff} in false acceptance and \plv{cmp.c1-c3.fr.diff} in false rejection
(Table~\ref{tab:comparisons}). In the discordant pairs, C1 accepts and C3 rejects \plv{cmp.c1-c3.fa.disc.ab} invalid
candidates, and C1 rejects and C3 accepts \plv{cmp.c1-c3.fr.disc.ba} valid ones. The pattern is the same with the audit
applied (C2 against C4: \plv{cmp.c2-c4.fa.diff} and \plv{cmp.c2-c4.fr.diff}).

Every invalid candidate that C1 accepts violates a requirement introduced by the correction (\plv{cv.c1} of the
correction violators). That follows largely from how the pool was built: \plv{neg.cv.n} of the \plv{labels.invalid}
invalid candidates violate a correction requirement, and \plv{neg.iv.n.w} does not. The \plv{pre.invalid.w} invalid
pre-correction states are accepted by C1 and C2 (\plv{pre.c1}, \plv{pre.c2}) and rejected by C3, C4, B2 and B1
(\plv{pre.c3}, \plv{pre.c4}, \plv{pre.b2}, \plv{pre.b1}). Those rejections are not blind: B2's generator and the
correction-aware contract proposer, whose contract C3 and C4 are generated from, saw the pre-correction state's edits,
abbreviated, in the session text (protocol~\S13; D41).

\paragraph{The audit changed an outcome.} Of the \plv{audit.verifiers.w} audited verifiers, \plv{audit.identical.w}
needed no revision and are identical to their plain counterparts (Table~\ref{tab:verifiers}). Across the
\plv{audit.outcomes} held-out outcomes of the \plv{audit.revised.w} revised verifiers, the audit changed
\plv{audit.changed.n.w}: C2 rejects an \texttt{inflection} near miss that C1 accepts. The audit cannot repair a verifier
whose contract is wrong. In \texttt{humanize} it flagged that the reference fails the initial-only tests, but a revision
that stays faithful to the contract cannot fix that (D27).

\paragraph{Baselines.} The session-LLM baseline (B2) accepts \plv{b2.fa} invalid candidates, a near miss that changes
the rounding of the default output, and rejects \plv{b2.fr} valid ones. Repository tests at the reference state (B1)
have the highest false-rejection rate (\plv{b1.fr}) and the lowest Youden's $J$ (\plv{b1.j}). Their pooled false
acceptance (\plv{b1.fa}) equals C2's (\plv{c2.fa}) and is below C1's (\plv{c1.fa}); their macro rate, which weights the
families equally, is slightly above C1's (\plv{b1.fa.macro} against \plv{c1.fa.macro}; Table~\ref{tab:summary}). These
caveats apply to B1:
\begin{itemize}
\item \emph{Agent-added tests.} The reference is the session agent's final state, so B1 includes
  \plv{b1.checks.added.min} to \plv{b1.checks.added.max} tests per family that the agent itself added
  (Table~\ref{tab:verifiers}). All \plv{b1.fr.added.w} of B1's false rejections come from those tests. Restricted to the
  base commit's own tests (exploratory), B1 rejects no valid candidate (\plv{b1.base.fr}) but accepts \plv{b1.base.fa}
  invalid ones. Both numbers largely follow from the label rule: a valid label requires passing the base commit's own
  tests (the regression norm), so this variant cannot reject a valid candidate, and the invalid candidates it rejects are
  those that fail the regression norm.
\item \emph{Not blind to construction.} The near-miss generators had B1's test files in their workspaces and ran them
  (D37), so B1's outcomes on near misses are not blind to construction.
\end{itemize}

\paragraph{The selection effect behind the zeros (exploratory).} D19 leaves as ambiguous the \plv{nm.ambiguous.w} near
misses whose violation the reference checks did not detect. These are the subtle cases:
\begin{itemize}
\item a stale file only after \texttt{os.chdir} with a relative path (\texttt{tinydb});
\item float positions rounded in the default grouping (\texttt{more-itertools});
\item plural acronyms such as ``URLs'' kept whole by \texttt{underscore} as well (\texttt{inflection});
\item acronyms joined by underscores (\texttt{inflection});
\item all-uppercase input with digits or symbols (\texttt{inflection}).
\end{itemize}
The correction-aware verifiers accept the \texttt{tinydb}, \texttt{more-itertools} and plural-acronym cases and reject
the other \texttt{inflection} cases. The reference checks and the
correction-aware verifiers are written by models of one model family of one provider, the checks from the reference
contract and the verifiers from the correction-aware contract, which covers every active reference requirement
(Table~\ref{tab:coverage}), and the writers of both saw the base repository and the reference diff, so their blind
spots are likely correlated. Split by label status, the correction-aware verifiers accept \plv{nm.conf.c3} of the
confirmed-invalid near misses but \plv{nm.amb.c3} of the ambiguous ones. If constructed candidates are labeled by
construction intent (Table~\ref{tab:sensitivity}), C3 and C4 accept \plv{intent.c3.fa} invalid candidates and B2
\plv{intent.b2.fa}; under the labels of the independent construction checks, C3 and C4 accept \plv{checks.c3.fa}. That
they accept no invalid candidate is therefore conditional on violations that the labeling checks can detect.

\IfFileExists{generated/results_comparisons.tex}{\begin{table}[t]\centering\footnotesize
\caption{Pre-specified paired comparisons on the same candidates, run \texttt{pilot-v0.2} (automated pilot). Last column: A accepts \& B rejects / A rejects \& B accepts.}\label{tab:comparisons}
\begin{tabular}{lllllrr}\toprule A & B & Metric & A$-$B & 95\% CI & n & discordant \\\midrule
C1 initial, plain & C3 full, plain & FA & +0.71 & [+0.50, +0.80] & 16 & 5/0 \\
C1 initial, plain & C3 full, plain & FR & +0.00 & [+0.00, +0.00] & 16 & 0/0 \\
C2 initial, audited & C4 full, audited & FA & +0.57 & [+0.00, +0.80] & 16 & 4/0 \\
C2 initial, audited & C4 full, audited & FR & +0.00 & [+0.00, +0.00] & 16 & 0/0 \\
C1 initial, plain & C2 initial, audited & FA & +0.14 & [+0.00, +0.50] & 16 & 1/0 \\
C1 initial, plain & C2 initial, audited & FR & +0.00 & [+0.00, +0.00] & 16 & 0/0 \\
C3 full, plain & C4 full, audited & FA & +0.00 & [+0.00, +0.00] & 16 & 0/0 \\
C3 full, plain & C4 full, audited & FR & +0.00 & [+0.00, +0.00] & 16 & 0/0 \\
B2 full-session LLM & C4 full, audited & FA & +0.00 & [+0.00, +0.00] & 16 & 0/0 \\
B2 full-session LLM & C4 full, audited & FR & +0.00 & [+0.00, +0.00] & 16 & 0/0 \\
B2 full-session LLM & C3 full, plain & FA & +0.00 & [+0.00, +0.00] & 16 & 0/0 \\
B2 full-session LLM & C3 full, plain & FR & +0.00 & [+0.00, +0.00] & 16 & 0/0 \\
B1 repository tests & C4 full, audited & FA & +0.43 & [+0.20, +1.00] & 16 & 3/0 \\
B1 repository tests & C4 full, audited & FR & +0.40 & [+0.33, +0.50] & 16 & 0/2 \\
\bottomrule\end{tabular}\end{table}
}{\unfilled{comparisons table}}
\IfFileExists{generated/sanity_table.tex}{\begin{table}[t]\centering\footnotesize
\caption{Sanity checks (excluded from the rates): reference solution and no-op passes per condition.}\label{tab:sanity}
\begin{tabular}{lrr}\toprule Condition & Reference passes & No-op passes \\\midrule
C1 initial, plain & 3/4 & 0/4 \\
C2 initial, audited & 3/4 & 0/4 \\
C3 full, plain & 4/4 & 0/4 \\
C4 full, audited & 4/4 & 0/4 \\
B2 full-session LLM & 4/4 & 0/4 \\
B1 repository tests & 4/4 & 0/4 \\
\bottomrule\end{tabular}\end{table}
}{\unfilled{sanity table}}
\IfFileExists{generated/sensitivity_table.tex}{\begin{table}[t]\centering\footnotesize
\caption{\textbf{Exploratory, not pre-specified:} rates when constructed candidates are labeled by construction intent (including those the pre-specified rule marked ambiguous).}\label{tab:sensitivity}
\begin{tabular}{lll}\toprule Condition & False acceptance & False rejection \\\midrule
C4 full, audited & 0.18 (2/11) [0.17, 0.20] & 0.00 (0/5) [0.00, 0.00] \\
C3 full, plain & 0.18 (2/11) [0.17, 0.20] & 0.00 (0/5) [0.00, 0.00] \\
C2 initial, audited & 0.64 (7/11) [0.40, 0.83] & 0.00 (0/5) [0.00, 0.00] \\
C1 initial, plain & 0.73 (8/11) [0.60, 0.83] & 0.00 (0/5) [0.00, 0.00] \\
B2 full-session LLM & 0.27 (3/11) [0.17, 0.40] & 0.00 (0/5) [0.00, 0.00] \\
B1 repository tests & 0.55 (6/11) [0.33, 0.80] & 0.40 (2/5) [0.33, 0.50] \\
\bottomrule\end{tabular}\end{table}
}{\unfilled{sensitivity table}}
\IfFileExists{generated/posthoc_table.tex}{\begin{table}[t]\centering\footnotesize
\caption{\textbf{Exploratory, post hoc (protocol v0.4).} Near misses accepted among those the label rule confirmed invalid and among those it marked ambiguous; rates under the independent construction checks' labels; rates without the two candidates that duplicate a development attack. Last row: B1 scored only on checks from the base commit's own tests (all candidates, primary labels).}\label{tab:posthoc}
\begin{tabular}{lrrrrrr}\toprule Condition & NM conf. & NM amb. & FA (checks) & FR (checks) & FA (no dup.) & FR (no dup.) \\\midrule
C1 initial, plain & 4/15 & 4/5 & 13/25 & 2/7 & 4/14 & 2/10 \\
C2 initial, audited & 3/15 & 4/5 & 12/25 & 2/7 & 3/14 & 2/10 \\
C3 full, plain & 0/15 & 3/5 & 6/25 & 0/7 & 0/14 & 0/10 \\
C4 full, audited & 0/15 & 3/5 & 6/25 & 0/7 & 0/14 & 0/10 \\
B2 full-session LLM & 1/15 & 4/5 & 8/25 & 0/7 & 1/14 & 0/10 \\
B1 repository tests & 5/15 & 3/5 & 10/25 & 3/7 & 4/14 & 4/10 \\
B1 upstream tests only & -- & -- & -- & -- & 13/17 (all) & 0/10 (all) \\
\bottomrule\end{tabular}\end{table}
}{\unfilled{post hoc table}}

\paragraph{Duplicated candidates.} The pool holds \plv{duplicates.w} held-out candidates that duplicate a development
attack used in their
family's C4 audit: the \texttt{more-itertools} pre-correction state and a \texttt{more-itertools} near miss that omits
the check of \texttt{step=0}, which have the attack's graded patch, and a \texttt{humanize} near miss that breaks custom formats,
which differs from its attack by a docstring line. Neither audit made a revision, but the pool counts the
\texttt{more-itertools} behaviour in both candidates that have its patch. Without the duplicates (exploratory), C1 accepts \plv{nodup.c1.fa}, C3 and C4
accept \plv{nodup.c3.fa}, and B1 accepts \plv{nodup.b1.fa}.

\paragraph{Label uncertainty.} The independent construction checks dispute the proposals of \plv{queue.disputed.w}
candidates that the rule labels valid:
\begin{itemize}
\item an \texttt{inflection} alternative, which diverges from the reference on exotic Unicode input and on some
  apostrophe-joined ASCII tokens;
\item the \texttt{inflection} pre-correction state, which diverges only on exotic Unicode input such as the Turkish
  dotted~İ;
\item a \texttt{tinydb} alternative, which fails read-after-write only when reopening the file fails.
\end{itemize}
C3 and C4 accept all of them. Together with the \plv{labels.ambiguous.w} ambiguous candidates, \plv{queue} of the
\plv{heldout} held-out candidates are queued for human review (Table~\ref{tab:labels}).

\IfFileExists{generated/labels_table.tex}{\begin{table}[t]\centering\footnotesize
\caption{Proposed labels of held-out candidates under rule D19 and the queue for human review (ambiguous, unlabeled, or disagreeing with the independent construction check).}\label{tab:labels}
\begin{tabular}{lrrrrr}\toprule Task & Valid & Invalid & Ambiguous & Unlabeled & Queued \\\midrule
CS-03 & 3 & 2 & 3 & 0 & 3 \\
CS-04 & 2 & 5 & 1 & 0 & 1 \\
\bottomrule\end{tabular}\end{table}
}{\unfilled{labels table}}

\paragraph{The reference checks' review.} Only CS-02's reference checks, written last, went through an independent
review-and-fix loop before freezing: at most \dv{pilot.refcheck.fixrounds.w} fix rounds, validated only on the
reference and the no-op (D33). The reference checks of CS-01, CS-03 and CS-04 had been frozen and used for label
proposals before that loop existed. They were not re-reviewed post hoc, because changing them after their labels were
seen would contaminate the labels (Section~\ref{sec:threats}).

\paragraph{Harness reliability.} A superseded run under harness 0.1.0 graded \plv{rel.v0.2.both} of the same pairs
(\plv{rel.v0.2.same} with an unchanged verifier and \plv{rel.v0.2.changed} more; Table~\ref{tab:reliability}). On the
pairs whose verifier did not change, its rewards agree with the final run in every case (\plv{rel.v0.2.reward}). The
other pairs are \texttt{more-itertools} B1 pairs graded with the earlier, truncated check selection (D31), and
\plv{rel.v0.2.changed.agree} of them agree; the differing pair is the pair that the harness resource-limit defects
turned into a \dv{pilot.stall.seconds}-second stall (D28--D30). The final run under harness \plv{run.harness} also
agrees, pair by pair, with the runs under 0.3.2, 0.3.1 and 0.3.0 (\plv{rel.v0.3.2.outcome}, \plv{rel.v0.3.1.outcome}
and \plv{rel.v0.3.outcome}). A re-grade of all \fv{rg.pairs} pairs under harness \fv{rg.harness}, each graded
\fv{rg.repeats} times (run \texttt{\REGRADERUN}, reliability only), gives the same outcome and the same reward on every
pair (\fv{rg.vs.pilotv033.outcome} and \fv{rg.vs.pilotv033.reward}; Table~\ref{tab:regrade}).

\IfFileExists{generated/reliability_table.tex}{\begin{table}[t]\centering\footnotesize
\caption{Agreement of superseded runs with the final run \texttt{pilot-v0.3} on pairs graded in both (pass vs non-pass; exact outcome). Superseded runs used earlier harness versions (D29--D31).}\label{tab:reliability}
\begin{tabular}{lrrr}\toprule Superseded run & Pairs & Reward agreement & Outcome agreement \\\midrule
\texttt{pilot-v0.2} & 108 & 108/108 & 108/108 \\
\texttt{pilot-v0.2-h2} & 4 & 4/4 & 4/4 \\
\bottomrule\end{tabular}\end{table}
}{\unfilled{reliability table}}
\IfFileExists{generated/fresh/regrade_reliability.tex}{\begin{table}[t]\centering\footnotesize
\caption{Reliability re-grade of the construction pilot (protocol \S15.3, \S16.1): run \texttt{pilot-v0.3.5-regrade} against the run it re-grades, on the pairs graded in both with an unchanged verifier (check-list hash). Pairs whose repeats disagree: n/a for a run graded once. Reward agreement is pass vs non-pass, outcome agreement the exact outcome. A reliability run, not a new pilot result.}\label{tab:regrade}
\begin{tabular}{lll}\toprule  & \texttt{pilot-v0.3.5-regrade} & \texttt{pilot-v0.3.3} \\\midrule
Harness & 0.3.5 & 0.3.3 \\
Repeats per pair & 3 & 1 \\
Graded on (UTC) & 2026-11-21 & 2026-11-18 \\
\quad from--to & 07:21:53--08:18:38 & 17:56:44--18:11:58 \\
Code revision & \texttt{a0542bb} & \texttt{c997779} \\
Pairs graded & 240 & 240 \\
Outcomes: pass/task failure/other & 112/128/0 & 112/128/0 \\
Pairs whose repeats disagree & 0 & n/a \\
Reward agreement with the re-grade & -- & 240/240 \\
Outcome agreement with the re-grade & -- & 240/240 \\
\bottomrule\end{tabular}\end{table}
}{\unfilled{re-grade table}}

\paragraph{Exploratory tests (D49).} A post hoc re-analysis gives exact sign-flip tests over the families
(Table~\ref{tab:expl-paired}). On false acceptance, C1 against C3 has one-sided $p=\plv{signflip.c1-c3.fa.p1}$ and
two-sided $p=\plv{signflip.c1-c3.fa.p2}$, at its conditional floor of \plv{signflip.c1-c3.fa.floor}, which CS-02, with no
difference, raises. As an equivalence question with a margin of \plv{tost.margin}, B2 against C3 and C4 cannot be shown
equivalent with \plv{run.families.w} families: the equivalence test's p never goes below \plv{tost.min}
(Table~\ref{tab:expl-equivalence}).

\IfFileExists{generated/exploratory/paired_signflip_table.tex}{\begin{table}[t]\centering\scriptsize
\caption{\textbf{Post hoc, exploratory (not pre-specified).} Paired comparisons of protocol \S7 with exact sign-flip randomization p-values over the 4 families (all sign assignments). Per-family and pooled differences A$-$B; discordant: A accepts \& B rejects / A rejects \& B accepts. $p_1$: one-sided, H1: A$-$B$>0$; $p_2$: two-sided; pooled (count-weighted) statistic, and $p_1$ for the family-mean statistic. Attainable floor: 0.0625 one-sided, 0.125 two-sided; cond.\ floor: the smallest one-sided p the families with zero difference allow.}\label{tab:expl-paired}
\begin{tabular}{llrrrrrrrrrr}\toprule A -- B & Metric & CS-01 & CS-02 & CS-03 & CS-04 & Pooled & Disc. & $p_1$ & $p_2$ & $p_1$ (mean) & Cond.\ floor \\\midrule
C1 -- C3 & FA & +0.17 & +0.00 & +0.50 & +0.80 & +0.35 & 6/0 & 0.125 & 0.25 & 0.125 & 0.125 \\
C1 -- C3 & FR & +1.00 & +0.00 & +0.00 & +0.00 & +0.20 & 0/2 & 0.5 & 1 & 0.5 & 0.5 \\
C2 -- C4 & FA & +0.17 & +0.00 & +0.00 & +0.80 & +0.29 & 5/0 & 0.25 & 0.5 & 0.25 & 0.25 \\
C2 -- C4 & FR & +1.00 & +0.00 & +0.00 & +0.00 & +0.20 & 0/2 & 0.5 & 1 & 0.5 & 0.5 \\
C1 -- C2 & FA & +0.00 & +0.00 & +0.50 & +0.00 & +0.06 & 1/0 & 0.5 & 1 & 0.5 & 0.5 \\
C1 -- C2 & FR & +0.00 & +0.00 & +0.00 & +0.00 & +0.00 & 0/0 & 1 & 1 & 1 & 1 \\
C3 -- C4 & FA & +0.00 & +0.00 & +0.00 & +0.00 & +0.00 & 0/0 & 1 & 1 & 1 & 1 \\
C3 -- C4 & FR & +0.00 & +0.00 & +0.00 & +0.00 & +0.00 & 0/0 & 1 & 1 & 1 & 1 \\
B2 -- C4 & FA & +0.17 & +0.00 & +0.00 & +0.00 & +0.06 & 1/0 & 0.5 & 1 & 0.5 & 0.5 \\
B2 -- C4 & FR & +0.00 & +0.00 & +0.00 & +0.00 & +0.00 & 0/0 & 1 & 1 & 1 & 1 \\
B2 -- C3 & FA & +0.17 & +0.00 & +0.00 & +0.00 & +0.06 & 1/0 & 0.5 & 1 & 0.5 & 0.5 \\
B2 -- C3 & FR & +0.00 & +0.00 & +0.00 & +0.00 & +0.00 & 0/0 & 1 & 1 & 1 & 1 \\
B1 -- C4 & FA & +0.33 & +0.00 & +1.00 & +0.20 & +0.29 & 5/0 & 0.125 & 0.25 & 0.125 & 0.125 \\
B1 -- C4 & FR & +0.00 & +1.00 & +0.33 & +0.00 & +0.40 & 0/4 & 0.25 & 0.5 & 0.25 & 0.25 \\
\bottomrule\end{tabular}\end{table}
}{\unfilled{exploratory sign-flip table}}
\IfFileExists{generated/exploratory/b2_equivalence_table.tex}{\begin{table}[t]\centering\scriptsize
\caption{\textbf{Post hoc, exploratory (not pre-specified).} B2 (full-session LLM tests) against C3/C4 as an equivalence question, margin $\pm$0.10: per-family and pooled differences, discordant pairs (B2 accepts \& C rejects / B2 rejects \& C accepts) and the TOST p-value (two one-sided sign-flip tests, pooled statistic). No p below 0.0625 is attainable, so equivalence cannot be concluded at one-sided $\alpha=0.05$. Labels: D19; the ambiguous near misses as invalid; the independent checks.}\label{tab:expl-equivalence}
\begin{tabular}{lllrrrrrrr}\toprule Pair & Labels & Metric & CS-01 & CS-02 & CS-03 & CS-04 & Pooled & Disc. & TOST $p$ \\\midrule
B2 -- C3 & D19 & FA & +0.17 & +0.00 & +0.00 & +0.00 & +0.06 & 1/0 & 0.25 \\
B2 -- C3 & D19 & FR & +0.00 & +0.00 & +0.00 & +0.00 & +0.00 & 0/0 & 0.0625 \\
B2 -- C3 & amb.\ as invalid & FA & +0.17 & +0.00 & +0.20 & +0.00 & +0.09 & 2/0 & 0.4375 \\
B2 -- C3 & amb.\ as invalid & FR & +0.00 & +0.00 & +0.00 & +0.00 & +0.00 & 0/0 & 0.0625 \\
B2 -- C3 & indep.\ checks & FA & +0.17 & +0.00 & +0.14 & +0.00 & +0.08 & 2/0 & 0.375 \\
B2 -- C3 & indep.\ checks & FR & +0.00 & +0.00 & +0.00 & +0.00 & +0.00 & 0/0 & 0.0625 \\
B2 -- C4 & D19 & FA & +0.17 & +0.00 & +0.00 & +0.00 & +0.06 & 1/0 & 0.25 \\
B2 -- C4 & D19 & FR & +0.00 & +0.00 & +0.00 & +0.00 & +0.00 & 0/0 & 0.0625 \\
B2 -- C4 & amb.\ as invalid & FA & +0.17 & +0.00 & +0.20 & +0.00 & +0.09 & 2/0 & 0.4375 \\
B2 -- C4 & amb.\ as invalid & FR & +0.00 & +0.00 & +0.00 & +0.00 & +0.00 & 0/0 & 0.0625 \\
B2 -- C4 & indep.\ checks & FA & +0.17 & +0.00 & +0.14 & +0.00 & +0.08 & 2/0 & 0.375 \\
B2 -- C4 & indep.\ checks & FR & +0.00 & +0.00 & +0.00 & +0.00 & +0.00 & 0/0 & 0.0625 \\
\bottomrule\end{tabular}\end{table}
}{\unfilled{exploratory equivalence table}}

\subsection{What the pilot's cases show}
\label{sec:analysis-pilot}
\paragraph{A retracted request inverts the reward (\texttt{humanize}).} The first request asked \texttt{naturalsize} to
drop a trailing \texttt{.0} by default. The correction kept the default output unchanged and made trimming opt-in. The
initial-only contract faithfully encodes the first request, and \plv{cov.cs01.initial.contradicting} of its
\plv{cov.cs01.initial.scored} scored requirements contradict the reference. The verifiers built from it therefore
reward the agent's own pre-correction state (\plv{pre.cs01.c1}), reject the \dv{pilot.pool.alternatives.w} independently written valid alternatives
(\plv{c1.cs01.fr}) and reject the reference solution itself (Table~\ref{tab:sanity}); Table~\ref{tab:matrix} gives every
outcome. This is the failure the problem statement predicts, the reward pointing at the behaviour the correction ruled
out, together with a failure it does not name: the reward also points away from the wanted behaviour. The audit detects the
inconsistency (the reference fails), but a revision that stays faithful to the contract cannot repair it.

\paragraph{An added requirement leaves edge cases open (\texttt{more-itertools}).} The first request added a
\texttt{step} argument. The correction added negative steps, an error for \texttt{step=0}, the positional
\texttt{ordering} argument and lazy groups, and kept the default output unchanged; \plv{cov.cs04.initial.missing} of these
requirements the initial-only contract lacks (negative steps, the error for \texttt{step=0}, laziness). The initial-only verifiers
accept the pre-correction state, where \texttt{step=0} silently returns groups (\plv{pre.cs04.c1}), and the near misses
that break negative steps, the error for \texttt{step=0} or laziness (\plv{nm.cs04.c1} of the invalid near misses). They reject
only the near miss that breaks positional ordering, which also breaks existing behaviour. Every correction-aware
condition rejects every invalid near miss and the pre-correction state (C3: \plv{nm.cs04.c3} and \plv{pre.cs04.c3}).

\paragraph{The first request can already imply the correction (\texttt{tinydb}).} For an atomic-write request, the
contract proposer inferred from the first message requirements covering everything the user said later
(covered: \plv{cov.cs02.initial.covered}; Table~\ref{tab:coverage}):
\begin{itemize}
\item reads through the same instance see the write;
\item the bytes are unchanged;
\item no temporary files are left behind;
\item the original file is intact after a failure.
\end{itemize}
These appear as \emph{inferred} requirements, which the acceptance gate requires a human to confirm, and initial-only
and correction-aware verifiers then make identical decisions. The information lost by initial-only reconstruction is a
property of the session \emph{and} the proposer, not of the session alone (Section~\ref{sec:threats}).

\paragraph{An initial-only contract can miss part of the correction (\texttt{inflection}).} The initial-only contract
lacks a correction requirement, all-uppercase input (\plv{cov.cs03.initial.missing} missing). The plain initial-only
verifier accepts the invalid near miss that breaks it, and the audited verifier rejects it (\plv{nm.cs03.c1} and
\plv{nm.cs03.c2} of the invalid near misses); both also accept \plv{nm.amb.cs03.c1} of the ambiguous near misses there.

\paragraph{The pre-correction state is not always invalid (\texttt{inflection}, \texttt{tinydb}).} In these sessions
the agent's state before the correction passes the reference checks and the regression norm, and the rule labels it
valid: the agent's state already met what the correction asked for. In \texttt{inflection}, for example, it already
handled ASCII all-uppercase input as before and left the helper functions' output unchanged. That case is informative,
but it means that the state before a scripted correction is not automatically an invalid candidate; some later
corrections restate what a careful agent would have done anyway. For \texttt{inflection} the independent check disputes
the label: it finds divergences from the reference limited to non-ASCII all-caps input. Whether ``behaves as before''
extends to the Turkish dotted~İ and circled capital letters is exactly the kind of judgement our labels must leave to
humans.

\paragraph{Repository tests fail in both directions.} Repository tests taken at the reference state, which include the
session agent's own tests:
\begin{itemize}
\item \emph{under-specify} the correction's invariants: they accept near misses that break custom formats or rounding in
  \texttt{humanize}, negative steps in \texttt{more-itertools}, and all-uppercase input and unchanged helper behaviour in
  \texttt{inflection};
\item and, through the tests the session agent added, \emph{over-specify} the reference's implementation: all
  \plv{b1.fr.added.w} false rejections come from those tests. They reject valid alternatives that lack a private reopen
  hook or the reference's clean-up of stale temporary files, and an alternative that splits CamelCase acronyms differently, and
  they also reject the valid \texttt{tinydb} pre-correction state.
\end{itemize}

\paragraph{The hard cases escape the verifiers (exploratory).} The near misses that the rule leaves ambiguous are those
whose violation appears only in a specific setup. The reference checks used for labeling miss them, so the rule cannot
confirm them, and the verifiers miss them too:
\begin{itemize}
\item \texttt{more-itertools}, floating-point positions in the default grouping: every condition accepts it;
\item \texttt{tinydb}, a relative database path followed by \texttt{os.chdir}: every generated-test condition accepts
  it, both correction-aware ones included;
\item \texttt{inflection}, plural acronyms: accepted by both correction-aware verifiers, B2 and B1, and rejected by both
  initial-only ones.
\end{itemize}
The other ambiguous \texttt{inflection} near misses (acronyms joined by underscores; all-uppercase input with digits or
symbols) are rejected by both correction-aware verifiers. If the construction claims are correct, leaving these
candidates out makes the rates optimistic about every condition. Settling that needs human labels and verifiers that
probe the setup space explicitly.

\IfFileExists{generated/outcome_matrix.tex}{\begin{table}[t]\centering\scriptsize
\caption{Outcome of every verifier on every held-out candidate, run \texttt{pilot-v0.2} (\checkmark{} pass, $\times$ any non-pass outcome).}\label{tab:matrix}
\begin{tabular}{lllcccccc}\toprule Candidate & Origin & Label & C4 & C3 & C2 & C1 & B2 & B1 \\\midrule
CS-03-eval-alt-1 & alternative & valid & \checkmark & \checkmark & \checkmark & \checkmark & \checkmark & \checkmark \\
CS-03-eval-alt-2 & alternative & valid & \checkmark & \checkmark & \checkmark & \checkmark & \checkmark & $\times$ \\
CS-03-eval-nm-1 & near miss & ambiguous & $\times$ & $\times$ & \checkmark & \checkmark & $\times$ & $\times$ \\
CS-03-eval-nm-2 & near miss & ambiguous & $\times$ & $\times$ & \checkmark & \checkmark & \checkmark & \checkmark \\
CS-03-eval-nm-3 & near miss & invalid & $\times$ & $\times$ & $\times$ & $\times$ & $\times$ & \checkmark \\
CS-03-eval-nm-4 & near miss & ambiguous & \checkmark & \checkmark & $\times$ & $\times$ & \checkmark & \checkmark \\
CS-03-eval-nm-5 & near miss & invalid & $\times$ & $\times$ & $\times$ & \checkmark & $\times$ & \checkmark \\
CS-03-pre-correction & pre-corr. & valid & \checkmark & \checkmark & \checkmark & \checkmark & \checkmark & \checkmark \\
CS-04-eval-alt-1 & alternative & valid & \checkmark & \checkmark & \checkmark & \checkmark & \checkmark & $\times$ \\
CS-04-eval-alt-2 & alternative & valid & \checkmark & \checkmark & \checkmark & \checkmark & \checkmark & \checkmark \\
CS-04-eval-nm-1 & near miss & invalid & $\times$ & $\times$ & \checkmark & \checkmark & $\times$ & \checkmark \\
CS-04-eval-nm-2 & near miss & invalid & $\times$ & $\times$ & \checkmark & \checkmark & $\times$ & $\times$ \\
CS-04-eval-nm-3 & near miss & invalid & $\times$ & $\times$ & $\times$ & $\times$ & $\times$ & $\times$ \\
CS-04-eval-nm-4 & near miss & invalid & $\times$ & $\times$ & \checkmark & \checkmark & $\times$ & $\times$ \\
CS-04-eval-nm-5 & near miss & ambiguous & \checkmark & \checkmark & \checkmark & \checkmark & \checkmark & \checkmark \\
CS-04-pre-correction & pre-corr. & invalid & $\times$ & $\times$ & \checkmark & \checkmark & $\times$ & $\times$ \\
\bottomrule\end{tabular}\end{table}
}{\unfilled{outcome matrix}}
